% !TeX program = xelatex
\documentclass[UTF8,12pt,a4paper]{ctexart}
\usepackage{amsmath,amssymb}
\usepackage{geometry}
\usepackage{tikz}
\geometry{margin=2cm}

% ===== 水印设置 =====
\usepackage{xcolor}
\usepackage{eso-pic}
\usepackage{graphicx}

\newcommand{\WMText}{贺昱琪学习资料}
\newcommand{\WMScale}{2.28}
\newcommand{\WMAngle}{45}
\colorlet{WMColor}{black!8}

\AddToShipoutPictureBG{%
  \AtPageLowerLeft{%
    \put(\LenToUnit{0.66\paperwidth},\LenToUnit{0.70\paperheight}){%
      \rotatebox{\WMAngle}{\scalebox{\WMScale}{\textcolor{WMColor}{\WMText}}}%
    }%
    \put(\LenToUnit{0.33\paperwidth},\LenToUnit{0.50\paperheight}){%
      \rotatebox{\WMAngle}{\scalebox{\WMScale}{\textcolor{WMColor}{\WMText}}}%
    }%
    \put(\LenToUnit{0.66\paperwidth},\LenToUnit{0.30\paperheight}){%
      \rotatebox{\WMAngle}{\scalebox{\WMScale}{\textcolor{WMColor}{\WMText}}}%
    }%
  }%
}

% ===== 8题铺满一页布局 =====
\usepackage{enumitem}
\newcommand{\BottomWriteSpace}{0.5cm}

\newenvironment{OnePageEightFill}{
  \small
  \vspace*{0pt}
  \begin{enumerate}[leftmargin=*, itemsep=0pt, topsep=0pt, parsep=0pt, partopsep=0pt]
  \setlength{\parskip}{0pt}
}{
  \end{enumerate}
  \vspace*{\BottomWriteSpace}
  \normalsize
  \newpage
}

\newcommand{\Q}[1]{\item #1\par\vfill}
\newcommand{\Qlast}[1]{\item #1\par}

\begin{document}

% =========================================================
% 8.3
% =========================================================
\section*{8月3日作业（8道题当晚22:00前打卡）}
\begin{OnePageEightFill}

\Q{解方程组：$\begin{cases} 2x - y = 7 \\ \dfrac{x-1}{2} = \dfrac{y+3}{3} \end{cases}$}

\Q{解方程组：$\begin{cases} 5x - y = 6 \\ \dfrac{x+4}{3} = \dfrac{y}{2} \end{cases}$}

\Q{（行程问题）甲、乙两人同时从 $A$ 地出发，以相同的速度向 $B$ 地前进。甲每行 5 分钟休息 1 分钟，乙每行 340 米休息 2 分钟。甲出发后 40 分钟到达 $B$ 地，乙到达 $B$ 地比甲迟了 2 分钟。已知两人最后一次的休息地点相距 140 米，求两人的速度均为每分钟多少米？}

\Q{（行程问题）海城地铁三号线呈西南至东北走向，全长约 78 千米。运行后，相关部门决定将地铁地下运行的平均速度提高 $\dfrac{1}{8}$，地上运行的平均速度降低 $\dfrac{1}{5}$，这样既可以让市民多欣赏 5 分钟风景，又能保证单程运行总时间不变。若地铁现在地下的平均速度比原来地上的平均速度每小时多行 4.5 千米。求：三号线单程运行多少分钟？三号线地上部分有多长（千米）？}

\Q{（钟表问题）从 6 点到 7 点，时针和分针成直角的时间是 6 时 \underline{\hspace{3em}} 分 和 6 时 \underline{\hspace{3em}} 分。}

\Q{（浓度问题）有 $A, B, C$ 三种盐水，按 $A$ 与 $B$ 数量之比为 $2:1$ 混合，得到浓度为 $13\%$ 的盐水；按 $A$ 与 $B$ 数量之比为 $1:2$ 混合，得到浓度为 $14\%$ 的盐水。如果 $A, B, C$ 数量之比为 $1:1:3$，混合成的盐水浓度为 $10.2\%$，问盐水 $C$ 的浓度是多少？}

\Q{（工程问题）有 $A, B$ 两个仓库，$A$ 仓库的货物是 $B$ 仓库的 2 倍。搬运完 $A$ 仓库的货物，甲单干需要 32 小时；搬运完 $B$ 仓库的货物，乙单干需要 24 小时，丙单干需要 12 小时。刚开始甲搬运 $A$ 仓库，乙搬运 $B$ 仓库，丙先帮甲，后来丙又去帮乙，直到最后两个仓库的货物同时搬完。则丙帮了甲几个小时？}

\Qlast{（行程问题）一列火车出发 1 小时后因故停车 0.5 小时，然后以原速的 $\dfrac{3}{4}$ 前进，最终到达目的地晚 1.5 小时。若出发 1 小时后又前进 90 公里再因故停车 0.5 小时，然后同样以原速的 $\dfrac{3}{4}$ 前进，则到达目的地仅晚 1 小时。那么整个路程为多少公里？}

\end{OnePageEightFill}

% =========================================================
% 8.4
% =========================================================
\section*{8月4日作业（8道题当晚22:00前打卡）}
\begin{OnePageEightFill}

\Q{解方程组：$\begin{cases} 3x + y = 17 \\ \dfrac{x+2}{3} = \dfrac{y-1}{2} \end{cases}$}

\Q{解方程组：$\begin{cases} 2x + 2y = 16 \\ \dfrac{x-1}{2} = \dfrac{y-2}{3} \end{cases}$}

\Q{（行程问题）甲、乙两人同时从 $A$ 地出发，以相同的速度向 $B$ 地前进。甲每行 4 分钟休息 1 分钟，乙每行 405 米休息 3 分钟。甲出发后 33 分钟到达 $B$ 地，乙到达 $B$ 地比甲迟了 3 分钟。已知两人最后一次的休息地点相距 225 米，求两人的速度均为每分钟多少米？}

\Q{（行程问题）某地铁线路全长约 100 千米。运行后，相关部门决定将地铁地下运行的平均速度提高 $\dfrac{1}{8}$，地上运行的平均速度降低 $\dfrac{1}{5}$，这样既可以让市民多欣赏 6 分钟风景，又能保证单程运行总时间不变。若地铁现在地下的平均速度比原来地上的平均速度每小时多行 20 千米。求：该地铁单程运行多少分钟？地上部分有多长（千米）？}

\Q{（钟表问题）从 7 点到 8 点，时针和分针成直角的时间是 7 时 \underline{\hspace{3em}} 分 和 7 时 \underline{\hspace{3em}} 分。}

\Q{（浓度问题）有 $A, B, C$ 三种盐水，按 $A$ 与 $B$ 数量之比为 $3:1$ 混合，得到浓度为 $12\%$ 的盐水；按 $A$ 与 $B$ 数量之比为 $1:3$ 混合，得到浓度为 $16\%$ 的盐水。如果 $A, B, C$ 数量之比为 $1:1:2$，混合成的盐水浓度为 $11\%$，问盐水 $C$ 的浓度是多少？}

\Q{（工程问题）有 $A, B$ 两个仓库，$A$ 仓库的货物是 $B$ 仓库的 3 倍。搬运完 $A$ 仓库的货物，甲单干需要 30 小时；搬运完 $B$ 仓库的货物，乙单干需要 20 小时，丙单干需要 10 小时。刚开始甲搬运 $A$ 仓库，乙搬运 $B$ 仓库，丙先帮甲，后来丙又去帮乙，直到最后两个仓库的货物同时搬完。则丙帮了甲几个小时？}

\Qlast{（行程问题）一列火车出发 2 小时后因故停车 0.5 小时，然后以原速的 $\dfrac{4}{5}$ 前进，最终到达目的地晚 1.5 小时。若出发 2 小时后又前进 120 公里再因故停车 0.5 小时，然后同样以原速的 $\dfrac{4}{5}$ 前进，则到达目的地仅晚 1 小时。那么整个路程为多少公里？}

\end{OnePageEightFill}

% =========================================================
% 8.5
% =========================================================
\section*{8月5日作业（8道题当晚22:00前打卡）}
\begin{OnePageEightFill}

\Q{解方程组：$\begin{cases} x + 3y = 12 \\ \dfrac{x-2}{2} = \dfrac{y+2}{2} \end{cases}$}

\Q{解方程组：$\begin{cases} x - y = 2 \\ \dfrac{x+1}{3} = \dfrac{y+1}{2} \end{cases}$}

\Q{（行程问题）甲、乙两人同时从 $A$ 地出发，以相同的速度向 $B$ 地前进。甲每行 6 分钟休息 2 分钟，乙每行 490 米休息 3 分钟。甲出发后 45 分钟到达 $B$ 地，乙到达 $B$ 地比甲迟了 2 分钟。已知两人最后一次的休息地点相距 140 米，求两人的速度均为每分钟多少米？}

\Q{（行程问题）某地铁线路全长约 74 千米。运行后，相关部门决定将地铁地下运行的平均速度提高 $\dfrac{1}{8}$，地上运行的平均速度降低 $\dfrac{1}{5}$，这样既可以让市民多欣赏 5 分钟风景，又能保证单程运行总时间不变。若地铁现在地下的平均速度比原来地上的平均速度每小时多行 21 千米。求：该地铁单程运行多少分钟？地上部分有多长（千米）？}

\Q{（钟表问题）从 8 点到 9 点，时针和分针成直角的时间是 8 时 \underline{\hspace{3em}} 分 和 \underline{\hspace{3em}} 时整（填分钟数或具体时刻）。}

\Q{（浓度问题）有 $A, B, C$ 三种盐水，按 $A$ 与 $B$ 数量之比为 $2:1$ 混合，得到浓度为 $14\%$ 的盐水；按 $A$ 与 $B$ 数量之比为 $1:2$ 混合，得到浓度为 $16\%$ 的盐水。如果 $A, B, C$ 数量之比为 $1:1:3$，混合成的盐水浓度为 $12\%$，问盐水 $C$ 的浓度是多少？}

\Q{（工程问题）有 $A, B$ 两个仓库，$A$ 仓库的货物是 $B$ 仓库的 2 倍。搬运完 $A$ 仓库的货物，甲单干需要 24 小时；搬运完 $B$ 仓库的货物，乙单干需要 15 小时，丙单干需要 10 小时。刚开始甲搬运 $A$ 仓库，乙搬运 $B$ 仓库，丙先帮甲，后来丙又去帮乙，直到最后两个仓库的货物同时搬完。则丙帮了甲几个小时？}

\Qlast{（行程问题）一列火车出发 1 小时后因故停车 1 小时，然后以原速的 $\dfrac{2}{3}$ 前进，最终到达目的地晚 3 小时。若出发 1 小时后又前进 80 公里再因故停车 1 小时，然后同样以原速的 $\dfrac{2}{3}$ 前进，则到达目的地仅晚 2 小时。那么整个路程为多少公里？}

\end{OnePageEightFill}

% =========================================================
% 8.6
% =========================================================
\section*{8月6日作业（8道题当晚22:00前打卡）}
\begin{OnePageEightFill}

\Q{解方程组：$\begin{cases} 2x - 3y = 2 \\ \dfrac{x-1}{3} = \dfrac{y}{2} \end{cases}$}

\Q{解方程组：$\begin{cases} 3x - 2y = 0 \\ \dfrac{x+2}{3} = \dfrac{y-2}{2} \end{cases}$}

\Q{（行程问题）甲、乙两人同时从 $A$ 地出发，以相同的速度向 $B$ 地前进。甲每行 7 分钟休息 3 分钟，乙每行 640 米休息 4 分钟。甲出发后 55 分钟到达 $B$ 地，乙到达 $B$ 地比甲迟了 1 分钟。已知两人最后一次的休息地点相距 240 米，求两人的速度均为每分钟多少米？}

\Q{（行程问题）某地铁线路全长约 63 千米。运行后，相关部门决定将地铁地下运行的平均速度提高 $\dfrac{1}{8}$，地上运行的平均速度降低 $\dfrac{1}{5}$，这样既可以让市民多欣赏 5 分钟风景，又能保证单程运行总时间不变。若地铁现在地下的平均速度比原来地上的平均速度每小时多行 27 千米。求：该地铁单程运行多少分钟？地上部分有多长（千米）？}

\Q{（钟表问题）从 5 点到 6 点，时针和分针成直角的时间是 5 时 \underline{\hspace{3em}} 分 和 5 时 \underline{\hspace{3em}} 分。}

\Q{（浓度问题）有 $A, B, C$ 三种盐水，按 $A$ 与 $B$ 数量之比为 $3:2$ 混合，得到浓度为 $13\%$ 的盐水；按 $A$ 与 $B$ 数量之比为 $2:3$ 混合，得到浓度为 $17\%$ 的盐水。如果 $A, B, C$ 数量之比为 $1:1:4$，混合成的盐水浓度为 $14\%$，问盐水 $C$ 的浓度是多少？}

\Q{（工程问题）有 $A, B$ 两个仓库，$A$ 仓库的货物是 $B$ 仓库的 3 倍。搬运完 $A$ 仓库的货物，甲单干需要 36 小时；搬运完 $B$ 仓库的货物，乙单干需要 18 小时，丙单干需要 9 小时。刚开始甲搬运 $A$ 仓库，乙搬运 $B$ 仓库，丙先帮甲，后来丙又去帮乙，直到最后两个仓库的货物同时搬完。则丙帮了甲几个小时？}

\Qlast{（行程问题）一列火车出发 1.5 小时后因故停车 0.5 小时，然后以原速的 $\dfrac{3}{4}$ 前进，最终到达目的地晚 2 小时。若出发 1.5 小时后又前进 90 公里再因故停车 0.5 小时，然后同样以原速的 $\dfrac{3}{4}$ 前进，则到达目的地仅晚 1.5 小时。那么整个路程为多少公里？}

\end{OnePageEightFill}

% =========================================================
% 8.7
% =========================================================
\section*{8月7日作业（8道题当晚22:00前打卡）}
\begin{OnePageEightFill}

\Q{解方程组：$\begin{cases} 4x + y = 19 \\ \dfrac{x+1}{2} = \dfrac{y-3}{2} \end{cases}$}

\Q{解方程组：$\begin{cases} x + y = 12 \\ \dfrac{x-3}{2} = \dfrac{y-1}{2} \end{cases}$}

\Q{（行程问题）甲、乙两人同时从 $A$ 地出发，以相同的速度向 $B$ 地前进。甲每行 5 分钟休息 2 分钟，乙每行 360 米休息 3 分钟。甲出发后 38 分钟到达 $B$ 地，乙到达 $B$ 地比甲迟了 8 分钟。已知两人最后一次的休息地点相距 90 米，求两人的速度均为每分钟多少米？}

\Q{（行程问题）某地铁线路全长约 68 千米。运行后，相关部门决定将地铁地下运行的平均速度提高 $\dfrac{1}{8}$，地上运行的平均速度降低 $\dfrac{1}{5}$，这样既可以让市民多欣赏 4 分钟风景，又能保证单程运行总时间不变。若地铁现在地下的平均速度比原来地上的平均速度每小时多行 15 千米。求：该地铁单程运行多少分钟？地上部分有多长（千米）？}

\Q{（钟表问题）从 4 点到 5 点，时针和分针成直角的时间是 4 时 \underline{\hspace{3em}} 分 和 4 时 \underline{\hspace{3em}} 分。}

\Q{（浓度问题）有 $A, B, C$ 三种盐水，按 $A$ 与 $B$ 数量之比为 $4:1$ 混合，得到浓度为 $12\%$ 的盐水；按 $A$ 与 $B$ 数量之比为 $1:4$ 混合，得到浓度为 $18\%$ 的盐水。如果 $A, B, C$ 数量之比为 $1:1:3$，混合成的盐水浓度为 $12\%$，问盐水 $C$ 的浓度是多少？}

\Q{（工程问题）有 $A, B$ 两个仓库，$A$ 仓库的货物是 $B$ 仓库的 2 倍。搬运完 $A$ 仓库的货物，甲单干需要 20 小时；搬运完 $B$ 仓库的货物，乙单干需要 30 小时，丙单干需要 15 小时。刚开始甲搬运 $A$ 仓库，乙搬运 $B$ 仓库，丙先帮甲，后来丙又去帮乙，直到最后两个仓库的货物同时搬完。则丙帮了甲几个小时？}

\Qlast{（行程问题）一列火车出发 2 小时后因故停车 1 小时，然后以原速的 $\dfrac{4}{5}$ 前进，最终到达目的地晚 2 小时。若出发 2 小时后又前进 160 公里再因故停车 1 小时，然后同样以原速的 $\dfrac{4}{5}$ 前进，则到达目的地仅晚 1.5 小时。那么整个路程为多少公里？}

\end{OnePageEightFill}

% =========================================================
% 8.8
% =========================================================
\section*{8月8日作业（8道题当晚22:00前打卡）}
\begin{OnePageEightFill}

\Q{解方程组：$\begin{cases} x - 4y = 4 \\ \dfrac{x-2}{6} = \dfrac{y+1}{2} \end{cases}$}

\Q{解方程组：$\begin{cases} 2x - y = 8 \\ \dfrac{x}{3} = \dfrac{y+2}{3} \end{cases}$}

\Q{（行程问题）甲、乙两人同时从 $A$ 地出发，以相同的速度向 $B$ 地前进。甲每行 8 分钟休息 2 分钟，乙每行 760 米休息 4 分钟。甲出发后 46 分钟到达 $B$ 地，乙到达 $B$ 地比甲迟了 8 分钟。已知两人最后一次的休息地点相距 160 米，求两人的速度均为每分钟多少米？}

\Q{（行程问题）某地铁线路全长约 136 千米。运行后，相关部门决定将地铁地下运行的平均速度提高 $\dfrac{1}{8}$，地上运行的平均速度降低 $\dfrac{1}{5}$，这样既可以让市民多欣赏 8 分钟风景，又能保证单程运行总时间不变。若地铁现在地下的平均速度比原来地上的平均速度每小时多行 15 千米。求：该地铁单程运行多少分钟？地上部分有多长（千米）？}

\Q{（钟表问题）从 10 点到 11 点，时针和分针成直角的时间是 10 时 \underline{\hspace{3em}} 分 和 10 时 \underline{\hspace{3em}} 分。}

\Q{（浓度问题）有 $A, B, C$ 三种盐水，按 $A$ 与 $B$ 数量之比为 $2:1$ 混合，得到浓度为 $10\%$ 的盐水；按 $A$ 与 $B$ 数量之比为 $1:2$ 混合，得到浓度为 $12\%$ 的盐水。如果 $A, B, C$ 数量之比为 $2:2:5$，混合成的盐水浓度为 $9\%$，问盐水 $C$ 的浓度是多少？}

\Q{（工程问题）有 $A, B$ 两个仓库，$A$ 仓库的货物是 $B$ 仓库的 3 倍。搬运完 $A$ 仓库的货物，甲单干需要 18 小时；搬运完 $B$ 仓库的货物，乙单干需要 24 小时，丙单干需要 8 小时。刚开始甲搬运 $A$ 仓库，乙搬运 $B$ 仓库，丙先帮甲，后来丙又去帮乙，直到最后两个仓库的货物同时搬完。则丙帮了甲几个小时？}

\Qlast{（行程问题）一列火车出发 3 小时后因故停车 1 小时，然后以原速的 $\dfrac{5}{6}$ 前进，最终到达目的地晚 2 小时。若出发 3 小时后又前进 250 公里再因故停车 1 小时，然后同样以原速的 $\dfrac{5}{6}$ 前进，则到达目的地仅晚 1.5 小时。那么整个路程为多少公里？}

\end{OnePageEightFill}

\end{document}